\documentclass[tikz,border=3mm,multi=tikzpicture]{standalone}
\usepackage{eleve-matematica-ef1}

\newcommand{\Termometro}[4]{%
  \draw[matematica contorno,line width=2pt] (#1,#2) circle (.48);
  \draw[matematica contorno,rounded corners=9pt,line width=2pt]
    (#1-.20,#2+.32) rectangle (#1+.20,#2+4.55);
  \fill[matemavermelho] (#1,#2) circle (.34);
  \fill[matemavermelho] (#1-.10,#2+.32) rectangle (#1+.10,#2+#3);
  \node[matematica resultado] at (#1,#2+5.20) {#4};
}

\begin{document}

% FIGURA 1 — fig-01-escala-do-termometro
\begin{tikzpicture}[matematica figura, x=1cm, y=1cm]
  \path[use as bounding box] (-0.30,-0.15) rectangle (5.90,7.40);
  \Termometro{2.00}{0.80}{2.40}{$20\,^\circ\mathrm{C}$}
  \foreach \y/\v in {1.12/0,2.08/10,3.04/20,4.00/30,4.96/40}{
    \draw[matematica eixo] (2.22,\y)--(2.65,\y);
    \node[font=\normalsize,text=matemacinza,anchor=west] at (2.82,\y) {\v};
  }
  \draw[matematica destaque,line width=2.2pt] (2.22,3.04)--(3.55,3.04);
\end{tikzpicture}

% FIGURA 2 — fig-02-maxima-e-minima
\begin{tikzpicture}[matematica figura, x=1cm, y=1cm]
  \path[use as bounding box] (-0.25,-0.15) rectangle (9.10,7.50);
  \Termometro{2.20}{0.80}{1.95}{$14\,^\circ\mathrm{C}$}
  \Termometro{6.70}{0.80}{3.55}{$27\,^\circ\mathrm{C}$}
  \node[matematica rotulo] at (2.20,0.05) {mínima};
  \node[matematica rotulo] at (6.70,0.05) {máxima};
\end{tikzpicture}

% FIGURA 3 — fig-03-grafico-de-temperaturas
\begin{tikzpicture}[matematica figura, x=1cm, y=.18cm]
  \path[use as bounding box] (-0.80,-4.00) rectangle (10.50,34.50);
  \draw[-{Stealth[length=3mm]},matematica eixo] (0,0)--(0,32);
  \draw[-{Stealth[length=3mm]},matematica eixo] (0,0)--(10,0);
  \foreach \y in {0,5,...,30}{
    \draw[matematica divisao] (-.18,\y)--(9.70,\y);
    \node[font=\normalsize,text=matemacinza,anchor=east] at (-.35,\y) {\y};
  }
  \foreach \x/\h/\d in {1.0/24/seg.,3.3/27/ter.,5.6/25/qua.,7.9/27/qui.}{
    \fill[matemazulclaro] (\x,0) rectangle ++(1.35,\h);
    \draw[matematica contorno] (\x,0) rectangle ++(1.35,\h);
    \node[font=\large\bfseries,text=matemazul] at (\x+.675,\h+1.6) {\h};
    \node[font=\normalsize,text=matemacinza] at (\x+.675,-2.0) {\d};
  }
\end{tikzpicture}

\end{document}
